% ------------------------------------------------------------------------
% file `1_E_juin-connaitre-2-exercise-body.tex'
%   in folder `xsim/'
%
%     exercise of type `connaitre' with id `2'
%
% generated by the `connaitre' environment of the
%   `xsim' package v0.21 (2022/02/12)
% from source `1_E_juin' on 2023/06/06 on line 202
% ------------------------------------------------------------------------
    Coche la ou les bonnes réponses.
    \begin{enumerate}
        \item Quels sont les nombres qui sont premiers ?
              \begin{qcm}
                  \item 12
                  \item 1
                  \item 7
              \end{qcm}
        \item Quel est l'élément caractéristique d'une translation ?
              \begin{qcm}
                  \item Le vecteur
                  \item L'axe de symétrie
                  \item Le centre de symétrie
              \end{qcm}
        \item Quels sont les nombres carrés ?
              \begin{qcm}
                  \item 30
                  \item 25
                  \item 50
              \end{qcm}
        \item Comment s'appelle une droite qui partage un segment en deux parts égales, perpendiculairement ?
              \begin{qcm}
                  \item Une médiane
                  \item Une médiatrice
                  \item Une bissectrice
              \end{qcm}
        \item Quel nom donne-t-on à un triangle ayant trois côtés de longueurs différentes ?
              \begin{qcm}
                  \item Isocèle
                  \item Équilatéral
                  \item Scalène
              \end{qcm}
        \item \(\widehat{A}\) et \(\widehat{B}\) sont deux angles supplémentaires. Quelle égalité traduit cette proposition ?
        \begin{qcm}
            \item \(\abs{\widehat{A}} + \abs{\widehat{B}} = \ang{180}\)
            \item \(\abs{\widehat{A}} +  \abs{\widehat{B}} = \ang{90}\)
            \item \(\abs{\widehat{A}} \cdot \abs{\widehat{B}} = \ang{180}\)
        \end{qcm}
        \item \(-3 \cdot 2 =\)
    \end{enumerate}
